% Part: incompleteness
% Chapter: incompleteness-provability
% Section: fixed-point-lemma

\documentclass[../../../include/open-logic-section]{subfiles}

\begin{document}

\olfileid[id]{inc}{inp}{fix}

\olsection{Lema Titik Tetap}

\begin{explain}
Lema titik tetap mengatakan bahwa untuk setiap !!{formula}~$!B(x)$, terdapat
!!a{sentence}~$!A$ sedemikian sehingga $\Th{T} \Proves !A \liff
!B(\gn{!A})$, asalkan $\Th{T}$ memperluas~$\Th{Q}$. Dalam kasus kalimat
pembohong, kita menginginkan $!A$ ekuivalen (secara terbukti
dalam~$\Th{T}$) dengan~``$\gn{!A}$ salah,'' yakni pernyataan bahwa
$\Gn{!A}$ adalah bilangan G\"odel dari !!{sentence} yang salah. Untuk
memahami gagasan pembuktiannya, akan berguna membandingkannya dengan uraian
informal Quine tentang~$!A$ sebagai ``{}`menghasilkan kepalsuan bila didahului
kutipannya sendiri' menghasilkan kepalsuan bila didahului kutipannya
sendiri.'' Operasi mengambil suatu ungkapan, lalu membentuk kalimat dengan
mendahului ungkapan itu dengan kutipannya sendiri dapat disebut
\emph{mendiagonalkan} ungkapan tersebut, dan hasilnya disebut
diagonalisasinya. Jadi, diagonalisasi `menghasilkan kepalsuan bila didahului
kutipannya sendiri' adalah ``{}`menghasilkan kepalsuan bila didahului
kutipannya sendiri' menghasilkan kepalsuan bila didahului kutipannya
sendiri.'' Perhatikan bahwa kalimat pembohong Quine bukan diagonalisasi
`menghasilkan kepalsuan,' melainkan diagonalisasi `menghasilkan kepalsuan
bila didahului kutipannya sendiri.' Jadi, sifat yang didiagonalkan untuk
menghasilkan kalimat pembohong itu sendiri melibatkan diagonalisasi!{}

Dalam bahasa aritmetika, kita membentuk kutipan !!a{formula} dengan satu
variabel bebas dengan menghitung bilangan G\"odel-nya, lalu menyubstitusikan
numeral standar bagi bilangan G\"odel tersebut ke variabel bebas.
Diagonalisasi~$!E(x)$ ialah $!E(\num{n})$, dengan $n = \Gn{!E(x)}$.
(Mulai sekarang, mari singkat $\num{\Gn{!E(x)}}$ sebagai $\gn{!E(x)}$.)
Jadi, jika $!B(x)$ adalah ``merupakan suatu kepalsuan,'' maka ``menghasilkan
kepalsuan bila didahului kutipannya sendiri'' menjadi ``menghasilkan
kepalsuan ketika diterapkan pada bilangan G\"odel dari diagonalisasinya.''
Jika kita memiliki simbol~$\Obj{diag}$ bagi fungsi $\fn{diag}(n)$ yang
menghitung bilangan G\"odel dari diagonalisasi !!{formula} dengan bilangan
G\"odel~$n$, kita dapat menuliskan $!E(x)$ sebagai
$!B(\Obj{diag}(x))$. Versi kalimat pembohong Quine kemudian menjadi
diagonalisasinya, yakni $!E(\gn{!E(x)})$ atau
$!B(\Obj{diag}(\gn{!B(\Obj{diag}(x))}))$. Tentu saja, $!B(x)$ kini dapat
berupa sifat lain apa pun, dan konstruksi yang sama akan bekerja. Untuk
teorema ketaklengkapan, kita ambil $!B(x)$ sebagai ``$x$~tidak
!!{derivable} dalam~$\Th{T}$.'' Maka $!E(x)$ berarti ``menghasilkan
!!a{sentence} yang tidak !!{derivable} dalam~$\Th{T}$ ketika diterapkan pada
bilangan G\"odel dari diagonalisasinya.''

Untuk memformalkan hal ini dalam~$\Th{T}$, kita harus menemukan cara untuk
memformalkan $\fn{diag}$. Fungsi $\fn{diag}(n)$ dapat dihitung; bahkan,
fungsi itu rekursif primitif: jika $n$ adalah bilangan G\"odel dari suatu
formula~$!E(x)$, $\fn{diag}(n)$ mengembalikan bilangan G\"odel
dari~$!E(\gn{!E(x)})$. (Ingat, $\gn{!E(x)}$ adalah numeral standar dari
bilangan G\"odel~$!E(x)$, yakni $\num{\Gn{!E(x)}}$.) Jika $\Obj{diag}$
merupakan simbol fungsi dalam $\Th{T}$ yang merepresentasikan fungsi
$\fn{diag}$, kita dapat mengambil $!A$ sebagai formula
$!B(\Obj{diag}(\gn{!B(\Obj{diag}(x))}))$. Perhatikan bahwa
\begin{align*}
\fn{diag}(\Gn{!B(\Obj{diag}(x))}) & =
\Gn{!B(\Obj{diag}(\gn{!B(\Obj{diag}(x))}))} \\
& = \Gn{!A}.
\end{align*}
Dengan mengasumsikan $\Th{T}$ dapat !!{derive}
\[
\Obj{diag}(\gn{!B(\Obj{diag}(x))}) = \gn{!A},
\]
teori itu dapat !!{derive} $!B(\Obj{diag}(\gn{!B(\Obj{diag}(x))}))
\liff !B(\gn{!A})$. Namun, ruas kiri menurut definisi adalah~$!A$.

Tentu saja, secara umum $\Obj{diag}$ bukan simbol fungsi~$\Th{T}$, dan jelas
bukan simbol fungsi~$\Th{Q}$. Namun, karena $\fn{diag}$ dapat dihitung,
fungsi itu \emph{dapat direpresentasikan} dalam~$\Th{Q}$ oleh suatu formula
$!D_{\fn{diag}}(x,y)$. Jadi, alih-alih menulis $!B(\Obj{diag}(x))$, kita
dapat menulis $\lexists[y][(!D_{\fn{diag}}(x,y) \land !B(y))]$. Selebihnya,
sketsa pembuktian di atas tetap berlaku, dan bahkan sudah berlaku
dalam~$\Th{Q}$.
\end{explain}

\begin{lem}
\ollabel{lem:fixed-point} Misalkan $!B(x)$ adalah sebarang formula dengan satu
variabel bebas~$x$. Maka terdapat suatu kalimat~$!A$ sedemikian sehingga
$\Th{Q} \Proves !A \liff !B(\gn{!A})$.
\end{lem}


\begin{proof}
Diberikan $!B(x)$, misalkan $!E(x)$ adalah formula
$\lexists[y][(!D_{\fn{diag}}(x,y) \land !B(y))]$ dan misalkan $!A$~adalah
diagonalisasinya, yakni formula $!E(\gn{!E(x)})$.

Karena $!D_{\fn{diag}}$ merepresentasikan $\fn{diag}$, dan
$\fn{diag}(\Gn{!E(x)}) = \Gn{!A}$, $\Th{Q}$ dapat !!{derive}
\begin{align}
  & !D_{\fn{diag}}(\gn{!E(x)}, \gn{!A}) \ollabel{repdiag1} \\
  & \lforall[y][(!D_{\fn{diag}}(\gn{!E(x)},y) \lif
  \eq[y][\gn{!A}])]. \ollabel{repdiag2}
\end{align}
Sekarang kita tunjukkan bahwa $\Th{Q} \Proves !A \liff !B(\gn{!A})$.
Kita menalar secara informal dengan hanya menggunakan logika dan fakta-fakta
yang !!{derivable} dalam~$\Th{Q}$.

Pertama, andaikan~$!A$, yakni $!E(\gn{!E(x)})$. Dengan kembali ke definisi
$!E(x)$, kita melihat bahwa $!E(\gn{!E(x)})$ tidak lain adalah
\[
\lexists[y][(!D_{\fn{diag}}(\gn{!E(x)},y) \land !B(y))].
\]
Perhatikan suatu~$y$ semacam itu. Karena
$!D_{\fn{diag}}(\gn{!E(x)},y)$, berdasarkan \olref{repdiag2},
$y = \gn{!A}$. Jadi, dari $!B(y)$ kita memperoleh~$!B(\gn{!A})$.

Sekarang andaikan $!B(\gn{!A})$. Berdasarkan \olref{repdiag1}, diperoleh
\begin{align*}
& !D_{\fn{diag}}(\gn{!E(x)}, \gn{!A}) \land !B(\gn{!A}).
\intertext{Maka}
& \lexists[y][(!D_{\fn{diag}}(\gn{!E(x)},y) \land !B(y))].
\end{align*}
Namun, itu tepat $!E(\gn{!E(x)})$, yakni~$!A$.
\end{proof}

\begin{digress}
Bandingkan pembuktian ini dengan pembuktian lema titik tetap dalam teori
komputabilitas. Perbedaannya ialah bahwa di sini kita hendak mendefinisikan
suatu \emph{pernyataan} dalam istilah dirinya sendiri, sedangkan di sana kita
hendak mendefinisikan suatu \emph{fungsi} dalam istilah dirinya sendiri;
terlepas dari perbedaan ini, gagasannya sungguh sama.
\end{digress}

\begin{prob}
  !!^a{formula}~$!A(x)$ adalah \emph{definisi kebenaran} jika $\Th{Q}
  \Proves !B \liff !A(\gn{!B})$ untuk semua !!{sentence}s~$!B$. Tunjukkan
  dengan menggunakan lema titik tetap bahwa tidak ada !!{formula} yang
  merupakan definisi kebenaran.
\end{prob}

\end{document}
